\documentclass[11pt]{article}

% Packages
\usepackage{geometry}
\usepackage{enumitem}
\usepackage{hyperref}
\usepackage{graphicx}
\usepackage{amsmath}
\usepackage{amssymb}
\usepackage{fancyhdr}
\usepackage[backend=biber,style=numeric-comp,sorting=none,date=year,url=false,isbn=false]{biblatex}
\addbibresource{background.bib}
% For arXiv-only entries, the DOI, URL and pre-published note repeat the eprint.
\AtEveryBibitem{\ifentrytype{online}{\iffieldundef{eprint}{}{%
  \clearfield{doi}\clearfield{url}\clearfield{urlyear}\clearfield{pubstate}}}{}}

% Page setup
\geometry{a4paper, margin=1in, headheight=13.6pt}
\pagestyle{fancy}
\fancyhf{}
\rhead{Plan Document}
\lhead{}
\cfoot{\thepage}

\title{Correlators of boosted resonances}
\author{Simon Rothman}
\date{\today}

\begin{document}

\maketitle

\section{Motivation}

Measuring EECs of vector bosons is highly appealing.
Potential applications include:
\begin{itemize}
    \item Search for boosted narrow resonances decaying hadronically (as a bump hunt in the EEC spectrum). (Are we sensitive to Higgs??).
    \item Calibration/proof-of-principle for EEC top mass measurement
    \item Measurement of transverse/longitudinal polarization fraction of high-pT Z and W
    \item Measurement of the full production spin-density matrix for high-pT Z and W bosons
\end{itemize}

\section{Background}

The two-point energy correlator (EEC) of a boosted, hadronically decaying boson looks very different from that of a QCD jet, whose collinear spectrum is smooth and approximately scale invariant~\cite{gaoPrecisionStudiesBoosted2026}.
A boson of boost \(\gamma = E_V/m_V\) instead produces a decay peak at \(z' \simeq 1/\gamma^2\), with \(z' = (1-\cos\theta')/2\), which at central rapidity corresponds to \(\Delta R \simeq 2m_V/p_T\)~\cite{gaoPrecisionStudiesBoosted2026}.
The width of the peak is set mainly by QCD radiation from the quarks rather than by \(\Gamma_V/m_V\), and well below the peak the spectrum follows that of a quark jet~\cite{ricciEnergyCorrelatorsHadronically2022}.

\subsection{Polarization and production}

In the narrow-width approximation, the correlator factorizes into a production density matrix and a decay matrix~\cite{ricciEnergyCorrelatorsHadronically2022,gaoPrecisionStudiesBoosted2026}.
The opening-angle spectrum depends on production only through the longitudinal fraction \(\rho_{00}\)~\cite{gaoPrecisionStudiesBoosted2026}.
Using NNLO production for \(pp \to V+X\) and decay correlators at NNLO with NNLL collinear and N\(^4\)LL back-to-back resummation, Gao et al.\ find that the peak height at \(\gamma = 4\) differs by about 70\% between purely longitudinal and purely transverse bosons.
For \(p_T > 100\)~GeV at \(y_V = 0\), the polarized spectrum is suppressed by up to about 10\% across the peak relative to an unpolarized treatment, because high-\(p_T\) production is dominantly transverse~\cite{gaoPrecisionStudiesBoosted2026}.
The off-diagonal entries require measuring the detector orientation relative to the production plane~\cite{gaoPrecisionStudiesBoosted2026}: \(|\Delta h| = 2\) interference appears in the azimuthal angle, while \(|\Delta h| = 1\) interference requires the azimuth and the angular separation jointly~\cite{ricciEnergyCorrelatorsHadronically2022}.
One-point energy and charge correlators together determine all eight real parameters of the density matrix~\cite{gaoPrecisionStudiesBoosted2026}.
Two experimental cautions come from Ref.~\cite{ricciEnergyCorrelatorsHadronically2022}: boosts along the LHC beam axis modify the azimuthal distributions, an effect expected to be mild but not yet assessed, and in a lepton-collider study \(D_2\) selections preferentially reject transverse bosons, homogenizing the longitudinal and transverse one-point spectra.

Two neighboring results extend the toolkit.
The one-point correlator of a vector current is fixed by a single parameter \(a_{\mathcal{E}}\), bounded by unitarity to \(-1/2 \leq a_{\mathcal{E}} \leq 1\); its N\(^3\)LO value at \(q^2 = m_Z^2\) is \(-0.42\), in good agreement with values extracted indirectly from LEP data~\cite{riembauFlowExtremalOnepoint2026}.
Flavor flux correlators generalize the program to flavor quantum numbers, with hadronization corrections to the \(b\) and \(c\) cases suppressed by \(\Lambda_{\mathrm{QCD}}/m_Q\)~\cite{leeFlavorFluxCorrelators2026}.

\subsection{Searches}

The same peak motivates resonance searches.
For a narrow resonance coupled linearly to a QCD operator, the peak position encodes the mass and the spectral shape encodes the spin and polarization, within a space of shapes bounded by unitarity and energy positivity~\cite{ricciBumpHuntingJets2026}.
A projection for a light hadrophilic \(Z'\) in \(pp \to Z'\gamma\), with statistical uncertainties only, finds that the inclusive EEC reaches sensitivity comparable to a recast of a CMS soft-drop mass search~\cite{ricciBumpHuntingJets2026}.
A separate recast of the CMS dijet correlator measurement also gives competitive limits, although its signal is at particle level and detector effects may degrade the reach~\cite{ricciBumpHuntingJets2026}.
In a Monte Carlo study of jets with \(p_T > 1\)~TeV recoiling against \(Z \to \ell\ell\), hadronic \(W/Z\) decays appear as a local excess of about 7\% even though they make up only about 1\% of the jets, with no reconstruction or tagging; that selection is better suited to the HL-LHC than to existing datasets~\cite{gaoPrecisionStudiesBoosted2026}.
The peak position also calibrates the angular scale, making \(m_V\) a standard candle inside jets, including the polarized \(W\) inside boosted top jets~\cite{gaoPrecisionStudiesBoosted2026}.


\section{Questions and Goals}

In the short term, we want to understand how we can use these ideas to do a CMS measurement. 
A few important lines of questioning:

\paragraph{``Untagged'' resonance search}
\begin{itemize}
    \item What sensitivity might be expect in an ``untagged''-style search for a narrow resonance in Run 3 data? 
    \item Which dataset(s) would we want to use? 
    \begin{itemize}
        \item Z+jet (tagged by Z$\to \ell\ell$)? Potential advantage: polarization of the other Z from leptonic decay?
        \item Photon+jet (tagged by photon)? Potential advantage: larger cross-section. Photon cannot be longitudinally polarized, I assume this has some implications for the polarization of the Z? 
        \item Dijet (with exactly two jets and momentum balance)? Potential advantage: large cross-section, but also large QCD bacground.
        \item QCD multijet (without momentum balance)? Potential advantage: largest cross-section, but also largest QCD background.
        \item Something else?
    \end{itemize}
    \item What do various SM processes look like?
    \begin{itemize}
        \item Can we separate W and Z at all?
        \item Can we see the Higgs? If not, can we set limits? Would they be competitive?
        \item What does the top look like? Do we learn anything interesting? 
    \end{itemize}
    \item Do we get any extra sensitivity from observables beyond two-point correlator? E.g,. E3C, E1C, etc. Including non-standard formulations thereof? The paper mentions E$^2$E$^2$ for example. Or with non-standard angles as in the one-point from the paper\ldots
    \item Would we get boosted sensitivity from some additional cuts? 
    \item Is heavy flavor interesting? Can we see heavy mesons? 
    \item Is the background really as clean as we hope it is? 
    \item How bad is the jet energy uncertainty? Is that the dominant systematic? What other systematics could be important? 
    \item Would our limits actually be competitive? 
\end{itemize}

\paragraph{Polarization measurements}

\begin{itemize}
    \item How good of a job do we need to do at tagging? Can we do it untagged? Is there any way to do good tagging that doesn't sculpt the distributions to shit? 
    \item Is there a useful shape difference between longitudinal and transverse bosons? Or is it dominantly the peak height? 
    \item Are the one-point correlators actually practical? What does the background do to them? 
    \item Would the three-point correlator be useful/interesting? Or some related things, maybe non-standard ones\ldots
    \item Is there an interesting bell inequality story here? Can we see interference effects? 
    \item Again, which dataset(s) would we want to use? 
    \item Charge correlators? Are they useful? How badly does hadronization etc.~ruin the picture? 
    \item How large are the spin analyzer coefficients?
    \item What is the other literature? Are we competitive? 
\end{itemize}

\section{Tasks and Timeline}
Break down the work into manageable tasks with deadlines.

\begin{enumerate}
    \item Task 1
    \item Task 2
    \item Task 3
\end{enumerate}

\section{Resources}
List required resources, dependencies, or tools.

\section{Risks and Mitigations}
Identify potential risks and how they will be addressed.

\section{Success Criteria}
Define what constitutes successful completion.

\printbibliography % chktex 1

\end{document}
